% ============================================================================
\section{Background: Payroll Income Tax in Indonesia}
\label{sec:background}
% ============================================================================
\subsection{Computing PPh 21 for a Permanent Employee}
\label{sec:pph21}
The monthly gross income of a permanent employee includes the salary, fixed allowances, irregular payments such as THR and bonuses, and benefits in kind, which have been taxable since July 2023. It also includes the premiums that the employer pays for work accidents, death, and health. Until 2023 the employer annualized each month's regular income and deducted the occupational expense, the employee's pension contributions, and the nontaxable amount (PTKP) of the employee's family status. The employer then taxed the result at the five progressive rates of 5\% to 35\% under Article 17~\cite{uu7_2021,pmk101_2016,per16_2016}. Since January 2024 the employer applies a monthly average effective rate (TER) to the gross income of January to November. The rate depends on the gross amount and on a category (A, B, or C) that follows from the PTKP status. In the final period the employer recomputes the tax of the whole year under Article 17 and withholds the difference from the amount already withheld~\cite{pp58_2023,pmk168_2023}. An employee with status K/1 and a monthly salary of Rp15,000,000 pays Rp933,660 in January 2024 at a TER of 6\% and Rp3,667,320 in March, when the THR doubles the gross income. The annual recomputation in December credits the earlier months and leaves Rp186,480 to withhold.

\subsection{Types and Dates of Rule Changes}
\label{sec:changes}
The corpus and the results are organized by the ten types of change of Table~\ref{tab:types}. Parameter changes (T1) are the most frequent among the real events, while company policies (T10) form the largest group of prospective units. Fig.~\ref{fig:timeline} places the 13 real regulations on a time axis with one color for each group of these types.

Two dates matter for every change. The issue date is the day on which the rule becomes known, whereas the start of validity is the day from which it binds. The start of validity of a tax facility such as DTP is the first tax period that it covers and can precede the date on which the regulation itself takes effect. The two dates coincide for most parameter changes. A retroactive rule separates them so that the correct withholding of a past month depends on the date on which the question is asked. The gray arrows of Fig.~\ref{fig:timeline} run from the issue date of each retroactive rule back to the start of its validity. The rounding rule of Director General Regulation PER-11/PJ/2025 is an example: it was issued in May 2025 for periods from January 2025. A KB that records only validity cannot reproduce what an employer should have withheld with the knowledge available at the time. A KB that records only the latest knowledge cannot report the correction owed for the months already paid.

\begin{table}[!t]
\caption{Types of change with their real events and prospective test units.}
\label{tab:types}
\centering\footnotesize
\setlength{\tabcolsep}{3pt}
\begin{tabularx}{\columnwidth}{@{}lLlr@{}}
\toprule
Code & Type of change & Real events & Units \\
\midrule
T1 & a single parameter value (wage cap, rate, ceiling) & E01, E05, E07, E08, E12 & 10 \\
T2 & a rate table is created or changed & E03 & 6 \\
T3 & income classification & E02, E04 & 6 \\
T4 & a new facility with eligibility criteria & E06, E11 & 7 \\
T5 & the scope of a facility & E10 & 6 \\
T6 & the validity period of a rule & E09 & 4 \\
T7 & a change of calculation method & E03, E04 & 3 \\
T8 & repeal without replacement & none & 3 \\
T9 & a calculation procedure such as rounding & E13 & 5 \\
T10 & a company policy & labor agreements & 30 \\
\bottomrule
\end{tabularx}
\end{table}

\subsection{Company Policies on Top of the Regulation}
\label{sec:policies}
Employers decide parts of the payroll themselves through company regulations or collective labor agreements. These parts include proration for partial months, the THR formula, overtime, allowances, and whether the employer bears the tax through a tax allowance (gross-up) or pays it directly. These policies form a second layer of rules that reads the values of the first and may not override mandatory provisions. A regulation can also apply differently across companies. PMK 10/2025 grants DTP only to employers whose business classification code is on its list. PP 7/2025 reduces the work-accident premium only for labor-intensive employers with at least 50 workers. The correctness of a payslip is therefore a property of each company because it depends on the regulation, the profile of the company, and the company's own policies taken together.

\subsection{Scope}
The study covers permanent employees of private and state-owned employers and tax years 2023--2026, with regulatory history from 2012 used for development. It excludes non-permanent workers, severance and other income taxed at final rates, and the filing of withholding slips.
